\section{Benchmarks}
\label{sec:benchmarks}

\subsection{How the numbers are read}
\label{sec:bench-method}

Every measurement below comes from a committed artefact in one of the
two repositories and is footnoted to it. We follow the measurement
contract those repositories impose on themselves
\citep{ecdlpconstant}, because its distinctions are exactly the ones
this paper needs to keep:
\begin{itemize}[leftmargin=1.4em,itemsep=0.2em]
\item \emph{Unit.} Where a whole pipeline is priced, the unit is
  $S=\text{operations}/\sqrt r$ in batched group additions, read
  against a matched signed-Frobenius Pollard rho on the same instance
  (Section~\ref{sec:rho-ref}) and the generic floor
  $\sqrt{\pi/(2\cdot2n)}$. A ratio over a smaller slice of the
  pipeline is a \emph{stage diagnostic} and is labelled as such.
\item \emph{Class.} A change is an \emph{advance} if the ratio to the
  floor fell, \emph{engineering} if $S$ fell and the ratio to the floor
  did not, \emph{relabelling} if a headline count fell while $S$ rose,
  and \emph{accounting} if only the unit changed. Nothing in this paper
  is an advance, and we say which of the other three each row is.
\item \emph{Isolation.} A CPU wall-clock ratio is promoted to a
  speedup only with a host-isolation receipt (exclusive cores, fixed
  frequency, A/A and quiet-host gates). Rows without one are
  \emph{exploratory}; deterministic counts (operations, rotations,
  comparisons, bytes, retired instructions) do not need one.
\item \emph{Correctness.} A row without a verified answer is not a
  result. ``Verified'' means every recovered logarithm was checked as
  $[\log]P_0=P$ against the planted scalar, in the run.
\end{itemize}

\subsection{Canonicalization cost}
\label{sec:bench-canon}

\paragraph{Normal-basis rotation against the squaring chain.}
Table~\ref{tab:canon-cost} prices one canonicalization at $n=61$ four
ways on one host: the obvious chain of $n-1$ squarings with a running
minimum (serial, and eight independent points in flight), the same
chain with shift reduction for a sparse irreducible, and the
normal-basis rotation of variant~A after a byte-table basis
change.\art{crypto:examples/koblitz_fold_cost.rs; crypto:research/notes/ecc2k130/RESEARCH_KOBLITZ_INDEX_CALCULUS.md, "The table was 2n times larger"}
In the same unit the sibling ledger prices a canonicalization at
$n=41$ at $50.8$~ns against $682.6$~ns for a batched group addition,
$0.074$ of an addition, and across the rho ladder of
Section~\ref{sec:bench-rho} at $0.09$--$0.12$ of an affine
addition.\art{crypto:research/notes/index-calculus/RESEARCH_IC_BOUNDARY_LEDGER.md, section 10.4 and section 19} For contrast, the Python model of the ECC2K-130 relation sweeps
measured a squaring-chain canonicalization at $398.52~\mu$s against
$29.20~\mu$s for a batched addition, $13.6$ additions, which is the
regime in which a folded table loses to a full
one.\art{crypto:research/notes/ecc2k130/RESEARCH_ECC2K130_RELATION_SWEEPS.md}

\begin{table}[t]
\centering\small
\caption{One canonicalization at $n=61$, nanoseconds per key, one host, exploratory. Variant~A is Section~\ref{sec:canon-impl}.}
\label{tab:canon-cost}
\begin{tabular}{@{}lrr@{}}
\toprule
method & ns / key & relative \\
\midrule
squaring chain, serial & $846.8$ & $9.9\times$ \\
squaring chain, 8 lanes & $733.7$ & $8.6\times$ \\
squaring with shift reduction & $1299.3$ & $15.2\times$ \\
normal-basis rotation (variant~A) & $\mathbf{85.5}$ & $1\times$ \\
\bottomrule
\end{tabular}
\end{table}

\paragraph{Zero gaps against the exhaustive scan.}
Table~\ref{tab:zerogap} is the frozen $n=83$ comparison of
variant~B with the exhaustive scan, on the same $2\times10^6$ sampled
pair states, the same $4\times10^6$ root records and the same one
target, six runs in the order baseline, zero-gap, zero-gap, baseline,
baseline, zero-gap.\art{experiments/koblitz-n83-index-design-20260929/packed_index_e1_canonical_20261009/README.md}
The rotation counts are deterministic: $656\,000\,000=8\times10^6\times82$
for the scan and $38\,261\,354$, or $4.78$ per word, for the shortcut, a
$17.1\times$ reduction. The canonical keys, shifts, first-witness
samples, exact-comparison counts and every non-canonicalization
operation count agreed between arms on all six runs. The timings are
on an unisolated host and are stage diagnostics: the three adjacent
target-stage pairs gave ratios $1.074$, $1.156$ and $1.118$, and the
canonicalization phase alone fell from $2.216$ to $1.557$~s
($1.42\times$). An adapter that replays exactly this stage through the
isolated benchmark service is committed; its receipt is not, so no
CPU speedup is claimed for the
shortcut.\art{experiments/koblitz-n83-index-design-20260929/packed_index_e1_isolated_replay_20261010/README.md}

\begin{table}[t]
\centering\small
\caption{Exhaustive scan against zero-gap canonicalization (variant~B) on the frozen $n=83$ probe. Counts are exact; times are medians of three on an unisolated host.}
\label{tab:zerogap}
\begin{tabular}{@{}lrr@{}}
\toprule
measure & exhaustive & zero-gap \\
\midrule
full rotations, index build & $328\,000\,000$ & $19\,127\,078$ \\
full rotations, target probe & $656\,000\,000$ & $38\,261\,354$ \\
exact key comparisons & $9\,324\,648$ & $9\,324\,648$ \\
root table, bytes & $88\,777\,220$ & $88\,777\,220$ \\
median index build & $2.130$~s & $1.980$~s \\
median conversion + canonicalization & $2.216$~s & $1.557$~s \\
median target stage & $7.204$~s & $6.511$~s \\
\bottomrule
\end{tabular}
\end{table}

\paragraph{Branches, lanes and vectors.}
Under Callgrind the portable key path, an exhaustive loop over eight
words, was $66.8\%$ of all instructions of the frozen $n=61$ cold
pipeline, about $586$ instructions per key over $74.4\times10^6$ keys.
A branch-free laned form of variant~A, measured on isolated cores
against its parent commit with identical keys, relations and
logarithms on all twelve cells, cut whole-pipeline time by
$1.09$--$1.31\times$ and retired instructions by $1.205\times$
($n=53$) and $1.398\times$ ($n=61$); at $n=41$ the gain was not
established. Class: engineering.\art{crypto:research/canon_laned_keys_20261008/README.md}
An exploration of the key kernels alone on two isolated cores put a
GFNI basis change at $0.82$--$0.94$~ns a key against $2.9$--$4.3$ for
the byte tables, and a funnel rotation at $2.4$--$2.8$~ns against
$5.4$--$7.1$ for the chained loop, with identical keys; this round is
declared and has no end-to-end result
yet.\art{crypto:research/ic_tool_program/explorations/R06-key-gfni-20261006/README.md}

\paragraph{Three limbs ($n=131$).}
At the ECC2K-130 degree a word takes three $64$-bit limbs ($64+64+3$
bits). A dependency-free crate implements five exact canonicalizers on
that layout and agrees with the exhaustive scan on every word for
$n\le20$ in several limb widths, on $10^7$ uniform $131$-bit words and
on structured families (low and high weight, zero runs straddling limb
boundaries, up to twelve tied longest runs, near-periodic words), with
four planted mutants all
caught.\art{experiments/canon-n131-20261014/README.md}
In retired instructions per word (cachegrind, uniform words), the
exhaustive scan costs $1850$, variant~B $1435$ ($5.43$ rotations a
word, but a ctz step for each of about $65$ one bits), a branch-free
refinement with a fixed $130$ steps and mask selects $4510$, and a
binary-lifting form of variant~A $514$, the cheapest and not
branch-free. On near-periodic words variant~B rises to $2738$, more
than the exhaustive scan, while both branch-free arms retire the same
count on every input; $72.4\%$ of uniform words have a single longest
run. No $131$-bit normal-basis multiply exists in the code base, so the
cost relative to a point operation was not measured, and the
nanosecond times of that report come from an unisolated host. The
branch-free form is therefore $2.4\times$ the exhaustive scan and
$8.8\times$ the lifting form at this degree: being free of branches
costs instructions here. Class: engineering, negative for the
branch-free arm.

\subsection{The canonicalization inside the reference}
\label{sec:bench-rho}

A signed-Frobenius rho walk that canonicalizes every step pays the
cost of Section~\ref{sec:bench-canon} once per step, and that cost
turns out to dominate the walk. Table~\ref{tab:rho-ladder} is the
instruction-counted ladder at $n=53$, $1024$ targets, in which each
rung contains the previous
ones.\art{crypto:research/notes/index-calculus/RESEARCH_STRONG_RHO_LADDER_20260929.md}
In the merged reference (rung R0) the canonicalization routine was
$50.82\%$ of retired instructions. From R1 (the same walk on the
library field) to R2, replacing it by the normal-basis least rotation
with the sign chosen by comparing $y$, together with a $\lambda^k$
table and a faster \texttt{mul\_mod}, cut the whole process
from $281.0$ to $68.9$ billion instructions, $4.08\times$, and still
left the (then $\Theta(n)$) least-rotation scan at $51.5\%$ of the
final rung. The index-calculus arm it was measured against moved from
$0.317$ (R1) to $1.294$ (R2) and $1.923$ (R3) times the reference's
cost without changing at all.
The compact-orbit panel of Section~\ref{sec:bench-e2e} shows the same
effect in CPU time: at $n=53$ with $1024$ targets the normal-basis
walk used $4.76$ CPU-seconds against $13.34$ for the same batched
engine with polynomial-basis canonicalization, $2.80\times$, while
performing $5\%$ \emph{more} steps.\art{crypto:research/notes/ecc2k130/compact_orbit_strong_rho_20260929/RESULT.md}

\begin{table}[t]
\centering\small
\caption{Retired instructions (Callgrind \texttt{Ir}) of the signed-Frobenius rho reference at $n=53$, $1024$ targets, all verified. R2 is the normal-basis canonicalization of Definition~\ref{def:sf}.}
\label{tab:rho-ladder}
\begin{tabular}{@{}llrr@{}}
\toprule
rung & change & Ir & IC / rho \\
\midrule
R0 & merged reference (software clmul) & $371\,102\,176\,689$ & $0.240$ \\
R1 & library field (hardware clmul) & $281\,027\,207\,019$ & $0.317$ \\
R2 & $+$ normal-coordinate canonicalization & $68\,926\,289\,783$ & $1.294$ \\
R3 & $+$ 32 lockstep lanes, batched inversion & $46\,384\,789\,571$ & $\mathbf{1.923}$ \\
\bottomrule
\end{tabular}
\end{table}

Two consequences follow. First, a comparison of an orbit-folded
index calculus against a rho that canonicalizes differently is a
comparison of two canonicalizations, not of two algorithms; the
``reversal'' in Section~\ref{sec:bench-e2e} is exactly that. Second,
Bailey et al.'s ECC2K-130 iteration $P\leftarrow P+\sigma^{j(P)}(P)$
with $j$ a function of the normal-basis Hamming weight is
class-invariant without canonicalizing each step, and normalizes only
distinguished points \citep{bailey2009}; the ledger prices a
canonical step at $2.74$--$2.92$ affine units against
$1.38$--$1.41$ for that step.\art{crypto:research/notes/index-calculus/RESEARCH_IC_BOUNDARY_LEDGER.md, section 19} The reference used here is therefore not the cheapest signed-Frobenius
rho that could be written, and every ``vs rho'' ratio below is
conservative in the index calculus's favour by up to that factor.

\subsection{Orbit compression of tables}
\label{sec:bench-tables}

\paragraph{The factor-base pair table, operation-counted ($n=41$).}
On the $39$-bit Koblitz rung of the boundary ledger the Frobenius fold
shrank the pair table from $6\,249\,999$ to $152\,439$ entries
($41.0\times$). Against the negation-folded base it replaced, the
factor-base phase fell from $8.45$ to $0.24$ in $S$ and relation
collection rose from $41.9$ to $45.0$, paid in canonicalizations;
net $S$ went from $50.4$ to $45.2$ (from $58.8$ for the unfolded
base, whose factor-base phase was $16.9$). On rungs with large cofactors, where the base was
almost the whole cost, the same fold is worth more: $n=29$ from
$34\,554$ to $728$, $n=31$ from $2489$ to $98.2$. Class:
engineering.\art{crypto:research/notes/index-calculus/RESEARCH_IC_BOUNDARY_LEDGER.md, section 10.4}

\paragraph{The folded tier at equal memory ($n=61$).}
Table~\ref{tab:fold-equal} compares the compact and folded tiers of
Section~\ref{sec:pairtable} at the same $0.64$~GiB on one host, $48$-bit
subgroup, two-summand descent.\art{crypto:docs/ic/runs/koblitz-orbit-fold-20260913.json}
At equal bytes the folded tier holds a base $11.1\times$ wider, which
divides the expected probes per target by $|F|^2$ and the probes
actually performed by $122\times$; the fold alone, on the code as it
stood before the row filter and orbit tag, was worth $2.0\times$, and
the rest of the $196\times$ per decomposed target is the wider base
made affordable. The prediction that a further $\sqrt2$ widening
should be worth about $1.9\times$ once recovery is $1.4\%$ of a
decomposition was tested on two runs and measured $1.93\times$ and
$2.04\times$, with every one of $21\,779$ and $22\,619$ targets
decomposed. These are wall-clock per-decomposition figures on an
unisolated host and a stage diagnostic; what they establish is that
the orbit tag reconnected the lever the fold exists for, not an
end-to-end speedup.

\begin{table}[t]
\centering\small
\caption{Compact against folded pair table at equal memory ($0.64$~GiB each), $n=61$, one host, exploratory. Hit rate is the fraction of residuals decomposed.}
\label{tab:fold-equal}
\begin{tabular}{@{}lrr@{}}
\toprule
 & compact & folded \\
\midrule
base points & $16\,592$ & $183\,488$ \\
stored pairs & $137\,655\,528$ & $138\,074\,720$ \\
probes per target & $1\,183\,373$ & $9\,676$ \\
ns per probe & $108$ & $331$ \\
hit rate & $0.0043$ & $0.9974$ \\
recovery per decomposition & $1.18$~ms & $0.026$~ms \\
seconds per decomposed target & $0.500$ & $0.00255$ \\
\bottomrule
\end{tabular}
\end{table}

\paragraph{What a probe costs, in the shape the pipeline probes.}
The fold's probe penalty depends on the access pattern: $2.8\times$
one target at a time, $2.0\times$ in blocks of $1024$ with prefetch,
$1.51\times$ inside the $m=3$ scan. In the scan the folded table costs
$74$~ns a base point more than the compact one, and the
canonicalization measured alone is $76$~ns: the fold's cost in the
descent is the canonicalization and nothing else, and blocking buys
the folded table $1.83\times$ because the $\sim76$~ns of dependent work
before the address is known is what stops consecutive probes'
memory round trips from
overlapping.\art{crypto:docs/ic/runs/koblitz-probe-shape-20260913.json}

\paragraph{The tier choice, priced in operations.}
When the default tier was flipped to fold-first on the wall-clock
evidence above and then re-priced over the whole cold pipeline in
batched additions, the flip did not survive: at $12\,688$ points
stored pairs fell $120.85\times$ while total operations rose $1.2\%$
(holdout targets: $1.016\times$). That is relabelling by definition.
Re-asked in a separate sweep over four widths with every other count
held identical, Table~\ref{tab:tier}, folded is $2.7\%$ dearer than
full at $12\,688$ points and $12.0\%$ cheaper at $18\,544$; each tier is cheapest somewhere: full below
about $8000$ points, compact to about $16\,000$, folded above, with
the full/folded crossover at $|F|\approx13\,623$ against a
pre-registered prediction of $13\,600$ from the measured per-probe
deltas. The twelve cells cost $12.1$--$15.4\times$ the matched rho's
$S$ and $18.9$--$24.1\times$ the floor on this instance; the tier moves the
constant only. Class: accounting for the sweep, relabelling for the
flip.\art{crypto:docs/ic/runs/koblitz-tier-crossover-20260921.json}

\begin{table}[t]
\centering\small
\caption{Total batched group additions of the cold pipeline at four factor-base widths, $n=61$, $47$-bit subgroup, 32 targets, all verified on every arm. Rho on the same instance: $S=0.1774$; floor $0.1135$.}
\label{tab:tier}
\begin{tabular}{@{}rrrrrl@{}}
\toprule
points & orbits & full & compact & folded & cheapest \\
\midrule
$6\,832$ & $56$ & $\mathbf{876\,342\,501}$ & $926\,956\,220$ & $1\,104\,202\,174$ & full \\
$9\,760$ & $80$ & $903\,649\,246$ & $\mathbf{880\,008\,741}$ & $1\,009\,804\,178$ & compact \\
$12\,688$ & $104$ & $941\,440\,872$ & $\mathbf{887\,500\,874}$ & $966\,559\,218$ & compact \\
$18\,544$ & $152$ & $1\,117\,432\,188$ & $1\,042\,518\,063$ & $\mathbf{982\,920\,004}$ & folded \\
\bottomrule
\end{tabular}
\end{table}

\paragraph{The packed root index ($n=83$).}
The record of Section~\ref{sec:record} holds the same $4\times10^6$
roots in $76.7$~MiB against $192.0$~MiB for the open-addressed hash
table ($2.50\times$), at a $1.41\times$ slower probe in the first
layout; the subsequent directory changes cut exact-key comparisons
from $58\,251\,024$ (raw prefix) to $25\,265\,242$ (hashed prefix),
$9\,324\,648$ ($22$-bit buckets) and $4\,943\,112$ (singleton tags),
all deterministic, with unisolated timings that moved in both
directions.\art{experiments/koblitz-n83-index-design-20260929/packed_index_e1_prefix_20261007/README.md, and the sibling E-1 directories}

\paragraph{Canonical pair tables as a cold job ($n=19,23$).}
Two preregistered, independently reviewed experiments in the research
ledger measured the orbit-canonical pair table against the expanded one
as complete fresh-process jobs.\art{crypto-autoresearcher:research/n19_pair_transport_20260921/REPORT.md; crypto-autoresearcher:research/pair_reuse_20260922/REPORT.md}
At $n=19$ (base of $152$ signed points, four orbit columns) the
canonical table has $293$ keys against $11\,097$, an $8.1\times$ smaller
key/value payload, and a paired whole-process ratio of $0.751$
$[0.655,0.864]$ (a $1.33\times$ speedup) over eight fixed targets; each
hit retains a real pair witness, the shift and sign are inverted, and
the recovered points are checked to sum to the residual. At $n=23$
with $16$ columns and $512$ queries per table build the ratio was
$0.647$ $[0.638,0.656]$, with $260\,361$ distinct pair sums held as
$5\,661$ canonical entries; at $n=19$ with $512$ queries the two were at
parity. The host was loaded; the Coordinator classified both as finite
implementation measurements and declined knowledge promotion.

\paragraph{SAT encodings.}
A representative-plus-shift encoding of the $n=53$ $S_5$ formula on
$220$ orbit columns has $22\,887$ clauses against $1\,815\,492$
($79.3\times$), pending independent
replay.\art{crypto:docs/ic/BOUNDARY_TARGETS.md, factor_base row} The same ledger found that an orbit-union encoding (one system per
orbit instead of one per shift) is cost-neutral for the solver,
ratios $1.000$--$1.467$ and never below $1$, and that a fixed-target
orbit constraint is unsound: Frobenius carries one target's system
to a conjugate target's system, not to
itself.\art{crypto-autoresearcher:knowledge/findings/KN-FIND-47da4e; crypto-autoresearcher:ledger/hypotheses/H-BINSTD-74b002.yaml} Compression of the formula is therefore not the same as a faster
solve, and the second claim is the one that has to be tested.

\subsection{End to end, with verified logarithms}
\label{sec:bench-e2e}

\paragraph{The $\Gamma$-invariant base with a Gr\"obner decomposition ($n=19$).}
Table~\ref{tab:n19} is the complete-pipeline comparison of
Section~\ref{sec:sage-pipeline} on $E_0/\FF_{2^{19}}$, ten planted
scalars, CPU seconds per ECDLP, speedup paired by scalar with a
bootstrap interval; all $40$ logarithms
verified.\art{experiments/volcano-ic/README.md, "E0 with a tau-invariant factor base"}
The orbit closure of a dimension-$6$ subspace needs $29$ unknowns
instead of $494$ and $39$ relations instead of $504$; it pays with
$19$ shifts per decomposition attempt ($5434$ solver calls against
$827$), and still finishes in $70$~s against $444$. The textbook
$\Gamma$-invariant base, normal-basis weight $\le3$, is a different
story: no $38$-variable Gr\"obner solve finished, and with a SAT solver
it is slower than the baseline. Three caveats are stated with the
result: the unit is CPU time on a shared machine, the timing windows
of the baseline were split by a hardware event, and the comparison is
against a plain subspace base, not against rho. This is the only
measurement in the paper in which a phase-aware base makes a complete,
verified index calculus faster, and it is a toy.

\begin{table}[t]
\centering\small
\caption{Complete index calculus on $E_0/\FF_{2^{19}}$, $r=130\,873$, ten scalars each, all logarithms verified. Speedup is the paired mean of baseline/variant CPU with a 95\% bootstrap interval.}
\label{tab:n19}
\resizebox{\textwidth}{!}{%
\begin{tabular}{@{}llrrrrl@{}}
\toprule
factor base & decomposition & unknowns & relations & solver calls & CPU s / ECDLP & vs baseline \\
\midrule
$x\in V$, $\dim10$ & Gr\"obner & $494$ & $504$ & $827$ & $444$ & $1\times$ \\
$\Gamma$-orbits of $x\in V'$, $\dim6$ & Gr\"obner, 19 shifts & $29$ & $39$ & $5434$ & $70$ & $\mathbf{6.85\times}$ $(5.70$--$8.02)$ \\
$\Gamma$-orbits of $x\in V'$, $\dim7$ & Gr\"obner, 19 shifts & $54$ & $64$ & $2664$ & $72$ & $6.26\times$ $(5.86$--$6.70)$ \\
normal-basis weight $\le3$ & CryptoMiniSat & $35$ & $45$ & $48$ & $668$ & $0.68\times$ $(0.62$--$0.74)$ \\
\bottomrule
\end{tabular}}
\end{table}

\paragraph{The folded pipeline against rho ($n=61$).}
The whole Rust pipeline with a $300\,608$-point, $2464$-column base
recovered $32$ of $32$ targets; the orbit tag and row filter cut
collection $31.4\times$, the solve $64.9\times$ and the descent
$3.8\times$ relative to the previous folded code, and peak memory from
$4.68$ to $2.40$~GiB. Its own record charges it $330.7\times$ the
matched rho with setup included ($86.7\times$ before these changes,
whose rho reference differed); the $12$--$15\times$ of
Table~\ref{tab:tier} belongs to the narrower $6.8$k--$18.5$k-point
sweep and is not this pipeline.\art{crypto:research/notes/ecc2k130/RESEARCH_KOBLITZ_INDEX_CALCULUS.md}

\paragraph{Compact-orbit index calculus against the strong reference ($n=37,41,53$).}
Table~\ref{tab:cold} is the isolation-gated cold panel: complete
fresh-process CPU of a compact-orbit index calculus (base
construction, full-rank collection, solve, verified recovery of one or
$1024$ targets) against the $32$-walk normal-basis signed-Frobenius
rho on the same public points, every cell past its preregistered A/A
and quiet-host gates, $135$ child processes, $46\,170$ target
logarithms agreed by an independent replay on a second
machine.\art{crypto:docs/ic/BOUNDARY_TARGETS.md, "Complete cold compact-orbit panel"; crypto:research/notes/ecc2k130/disjoint_cold_v2_outcome_20261001/RESULT.md}
Every interval lies above parity. An earlier, separate experiment on
different public points (2026-09-29, not isolation-gated) measured both
references in one run: at its three batch cells the index calculus
read $0.78$, $0.39$ and $0.34$ of the \emph{previous} rho reference
(PCLMUL, Itoh--Tsujii), an apparent index-calculus win, and $2.08$,
$2.99$ and $3.09$ of the normal-basis reference. Those six numbers
come from that experiment and are not cells of Table~\ref{tab:cold},
which reads $3.03$, $1.31$ and $1.61$ at the same nominal cells. The
win was the reference's canonicalization, which is the point of
Section~\ref{sec:bench-rho}.\art{crypto:research/notes/ecc2k130/compact_orbit_strong_rho_20260929/RESULT.md}

\begin{table}[t]
\centering\small
\caption{Complete cold CPU, compact-orbit index calculus over signed-Frobenius rho, paired median with 95\% interval, isolation-gated. $K$ is the number of orbit columns.}
\label{tab:cold}
\begin{tabular}{@{}lllr@{}}
\toprule
cell & policy & IC / rho & reduction to parity \\
\midrule
$n=37$, 1 target, $K=7$ & unfiltered & $2.810$ $[2.718,2.874]$ & $64.4\%$ \\
$n=37$, 1024 targets, $K=42$ & unfiltered & $3.031$ $[2.998,3.047]$ & $67.0\%$ \\
$n=41$, 1 target, $K=85$ & unfiltered & $13.378$ $[9.367,18.491]$ & $92.5\%$ \\
$n=41$, 1024 targets, $K=255$ & blocked prefilter & $1.314$ $[1.288,1.359]$ & $23.9\%$ \\
$n=53$, 1 target, $K=220$ & unfiltered & $22.675$ $[17.904,31.595]$ & $95.6\%$ \\
$n=53$, 1024 targets, $K=440$ & blocked prefilter & $1.605$ $[1.566,1.666]$ & $37.7\%$ \\
\bottomrule
\end{tabular}
\end{table}

\paragraph{Online against cold: one target ($n=41$--$83$).}
Single-target comparisons read differently depending on whether the
index calculus's setup (base, relations, logarithm table) is charged.
The boundary ledger's isolated single-target round separates the two
readings on $64$ public points at each of six sizes,
$r=2^{36.6}$--$2^{47.2}$, both arms on the same
point.\art{crypto:research/notes/index-calculus/RESEARCH_IC_BOUNDARY_LEDGER.md, section 23}
With setup excluded the index calculus's online interval is
$8.8$--$17.5\times$ faster than single-target rho, every interval
above $1$ ($4.3$--$9.4\times$ if rho is priced at the canonical
step). With setup included it is $2.0$--$31.6\times$ slower at every
size, and the setup pays for itself only after $7$--$37$ targets. The
same setup used as precomputation for Bernstein and Lange's generic
walk would, by the ledger's model, reach an online cost $2.0$--$8.0\times$
below the index calculus's, so the online advantage belongs to the
precomputation and not to the method. A separate isolated panel at
$n=61$, $a=0$ ($64$ unseen points per host, two hosts) reads
$12.9\times$ and $13.9\times$ online ($95\%$ interval $9.8$--$19.1$ on
the pinned host) with the setup excluded and reported as a cold ratio
of $83.4$.\art{crypto:docs/ic/BOUNDARY_TARGETS.md, vs_rho row} The
earlier $n=83$, $a=1$ observation (three runs on one frozen point,
median online ratio $8.2$, shared host, no isolation receipt) is of
the same online kind, and the ledger now keeps it as
history.\art{crypto:experiments/koblitz-single-target-n83-20261006/}
It was also measured against a packed rho build that the repository no
longer accepts as a reference, whose jump setup depends on the target.
Five interleaved pairs on the same point against the current strong
signed-Frobenius rho, on one virtualized host and labelled diagnostic,
read index calculus over rho online at $1.38$ ($95\%$ interval
$1.04$--$1.88$): rho was faster. A follow-up against a tuned rho on
five fresh targets found no uniform online winner, its interval
containing parity.\art{crypto:research/n83_evaluation_20261012/README.md; crypto:research/n83_evaluation_20261012/optimization/README.md}
None of these is an advance under Section~\ref{sec:bench-method}: each
reads a method with precomputation against a reference without it.

\subsection{What survives complete relation collection}
\label{sec:survives}

Table~\ref{tab:survives} collects the paper's improvements with what
each one moved and its class. The pattern is uniform. Representation
changes move constants, sometimes large ones, inside a method that
remains slower than the matched rho at every degree measured once its
setup is charged ($1.3$--$331\times$ across the rungs priced) and is
faster online only by spending precomputation that a generic walk
could spend better (Section~\ref{sec:bench-e2e}); the same representation change inside the reference moves
its constant by a comparable factor; and the headline count that a
fold divides by $2n$ (entries, keys, clauses, unknowns) is never the
count the end-to-end cost is proportional to.

\begin{table}[t]
\centering\small
\caption{Each improvement, what it changed, and its class under the contract of Section~\ref{sec:bench-method}.}
\label{tab:survives}
\begin{tabularx}{\textwidth}{@{}>{\raggedright\arraybackslash}p{0.26\textwidth}X>{\raggedright\arraybackslash}p{0.13\textwidth}>{\raggedright\arraybackslash}p{0.12\textwidth}@{}}
\toprule
improvement & what moved & verified & class \\
\midrule
normal-basis rotation vs squaring chain & $8.6$--$15\times$ per key; rho process $4.08\times$ fewer instructions & keys identical; 1024/1024 & engineering \\
zero-gap shortcut (variant B) & $17.1\times$ fewer rotations; canonicalization phase $1.42\times$ (unisolated) & keys and shifts identical & engineering (stage) \\
laned variant A & $1.09$--$1.31\times$ pipeline (isolated) & identical logs & engineering \\
packed 18-byte record & $2.27\times$ fewer table bytes ($2.50\times$ first layout); $11.8\times$ fewer comparisons across successive layouts on one frozen probe & controls identical & engineering (stage) \\
folded pair table & $2n$ fewer entries; $+1.2\%$ operations at 12688 pts; cheapest above $\approx16\,000$ pts & 32/32 & relabelling / accounting \\
orbit tag & recovery $O(|F|)\to O(n)$; widening lever restored ($1.9\times$ predicted, $1.93$--$2.04$ measured) & all targets decomposed & engineering \\
canonical pair table as a cold job & $0.75$ ($n=19$) and $0.65$ ($n=23$) whole-job time, two separate experiments & witnesses verified & engineering \\
$\Gamma$-orbit Gr\"obner base ($n=19$) & $6.85\times$ CPU per complete ECDLP vs plain base & 40/40 & engineering \\
representative+shift SAT encoding & $79.3\times$ fewer clauses; solve time unchanged & pending replay & relabelling \\
\bottomrule
\end{tabularx}
\end{table}
